\documentclass[aspectratio=169,12pt]{beamer}
% Maths Lycée — thème éditorial inspiré du portfolio de Yassine.
% Modifier ici les couleurs et la typographie de tous les diaporamas.
\usepackage[T1]{fontenc}
\usepackage[utf8]{inputenc}
\usepackage[french]{babel}
\usepackage{lmodern,amsmath,amssymb,mathtools,tikz,pgfplots,listings,adjustbox}
\usetikzlibrary{arrows.meta,positioning,calc}
\pgfplotsset{compat=1.18,width=9cm,height=4.2cm}
\definecolor{Ink}{HTML}{202124}
\definecolor{Blue}{HTML}{356B9E}
\definecolor{Mist}{HTML}{F6F6F5}
\definecolor{Muted}{HTML}{737579}
\definecolor{Line}{HTML}{E7E7E5}
\setbeamercolor{normal text}{fg=Ink,bg=white}
\setbeamercolor{structure}{fg=Blue}
\setbeamercolor{frametitle}{fg=Ink}
\setbeamercolor{title}{fg=Ink}
\setbeamercolor{block title}{fg=Blue,bg=Mist}
\setbeamercolor{block body}{fg=Ink,bg=Mist}
\setbeamerfont{title}{size=\LARGE,series=\bfseries}
\setbeamerfont{frametitle}{size=\Large,series=\bfseries}
\setbeamersize{text margin left=1cm,text margin right=1cm}
\setbeamertemplate{navigation symbols}{}
\setbeamertemplate{itemize items}[circle]
\setbeamertemplate{enumerate items}[default]
\setlength{\parskip}{7pt}
\newcommand{\slidekind}{COMPRENDRE}
\newcommand{\chapternumber}{01}
\setbeamertemplate{frametitle}{\vspace{0.3cm}{\fontsize{7}{9}\selectfont\color{Blue}\MakeUppercase{\slidekind}}\par\vspace{0.1cm}\insertframetitle\par\vspace{0.08cm}\color{Line}\rule{\textwidth}{0.4pt}\vspace{-0.02cm}}
\setbeamertemplate{footline}{\hspace{1cm}\begin{minipage}{\dimexpr\paperwidth-2cm}\color{Line}\rule{\textwidth}{0.4pt}\par\vspace{0.06cm}{\fontsize{6.5}{8}\selectfont\color{Muted}MATHS LYCÉE\quad / \quad PREMIÈRE SPÉCIALITÉ\hfill\insertframenumber\ /\ \inserttotalframenumber}\end{minipage}\vspace{0.22cm}}
\setbeamertemplate{title page}{
\vspace{0.3cm}
{\fontsize{8}{10}\selectfont\color{Blue}MATHS LYCÉE\quad /\quad 2026--2027}\par
\vspace{0.6cm}
\begin{columns}[c,onlytextwidth]
\begin{column}{.72\textwidth}
{\fontsize{9}{12}\selectfont\color{Muted}CHAPITRE \chapternumber}\par\vspace{0.22cm}
{\usebeamerfont{title}\color{Ink}\inserttitle\par}
\vspace{0.4cm}{\color{Blue}\rule{1.5cm}{2pt}}\par
{\small\color{Muted}Comprendre les notions.\par Justifier chaque étape.}
\end{column}
\begin{column}{.23\textwidth}
\begin{tikzpicture}[scale=.65]
\fill[Mist,rounded corners=6pt](-1.9,-2.1)rectangle(2,2.5);
\draw[Line,step=.5](-1.8,-2)grid(1.9,2.4);
\draw[Muted,->](-1.7,-1.3)--(1.8,-1.3);
\draw[Muted,->](0,-1.6)--(0,2.2);
\draw[Blue,very thick,domain=-1.5:1.5,samples=40]plot(\x,{1.3*\x*\x-1.3});
\fill[Blue](1,0)circle(2.5pt);
\end{tikzpicture}
\end{column}
\end{columns}
\vfill}
\lstdefinestyle{mathspython}{language=Python,basicstyle=\ttfamily\small,keywordstyle=\color{Blue}\bfseries,commentstyle=\color{Muted},stringstyle=\color{Blue},backgroundcolor=\color{Mist},frame=none,xleftmargin=8pt,xrightmargin=8pt,aboveskip=8pt,belowskip=8pt,showstringspaces=false,columns=fullflexible,keepspaces=true,breaklines=true}
\lstnewenvironment{codepython}{\lstset{style=mathspython}}{}
\newcommand{\R}{\mathbb{R}}
\newcommand{\N}{\mathbb{N}}
\newcommand{\E}{\mathbb{E}}
\newcommand{\V}{\mathrm{V}}
\hypersetup{colorlinks=true,urlcolor=Blue}

\title{Droites et cercles dans un repère}
\subtitle{Première spécialité}
\renewcommand{\chapternumber}{10}\author{Maths Lycée}
\date{Année scolaire 2026--2027}
\begin{document}\begin{frame}\titlepage\end{frame}
\begin{frame}{Objectifs du chapitre}\begin{itemize}
\item Lire et construire un vecteur normal
\item Déterminer une équation de droite ou de cercle
\item Calculer un projeté orthogonal
\item Résoudre un problème d’intersection
\end{itemize}\end{frame}
\renewcommand{\slidekind}{COMPRENDRE}
\begin{frame}[t,fragile]{Un problème de localisation}
Une antenne est située en $C(2,1)$ et sa portée est de $5$ unités. Les points exactement à la limite de portée vérifient $CM=5$.\medskip\par Une route suit la droite $y=1$. Pour savoir où elle coupe cette limite, on peut traduire les deux conditions en équations. Tout le chapitre relie ainsi figures et calculs.
\end{frame}
\renewcommand{\slidekind}{COMPRENDRE}
\begin{frame}[t,fragile]{Le cadre : un repère orthonormé}
Les axes sont perpendiculaires et utilisent la même unité. Pour $A(x_A,y_A)$ et $B(x_B,y_B)$ :\[\overrightarrow{AB}=(x_B-x_A,y_B-y_A),\]\[AB^2=(x_B-x_A)^2+(y_B-y_A)^2.\]Le milieu de $[AB]$ a pour coordonnées $\left(\frac{x_A+x_B}{2},\frac{y_A+y_B}{2}\right)$.
\end{frame}
\renewcommand{\slidekind}{COMPRENDRE}
\begin{frame}[t,fragile]{Vecteur normal à une droite}
Un vecteur \textbf{normal} à une droite est un vecteur non nul orthogonal à un vecteur directeur de cette droite.\medskip\par Si $\vec n=(a,b)$ est normal, alors $\vec d=(-b,a)$ dirige la droite, car $a(-b)+ba=0$.\medskip\par Exemple : $(2,3)$ est normal à toute droite de direction $(-3,2)$.
\end{frame}
\renewcommand{\slidekind}{MÉTHODE \& EXEMPLE}
\begin{frame}[t,fragile]{D’une condition géométrique à une équation}
La droite passant par $A(x_A,y_A)$ et de normale $\vec n=(a,b)\ne(0,0)$ est définie par\[\vec n\cdot\overrightarrow{AM}=0.\]Pour $M(x,y)$, cela donne\[a(x-x_A)+b(y-y_A)=0.\]En développant, on obtient $ax+by+c=0$, où $c=-ax_A-by_A$.
\end{frame}
\renewcommand{\slidekind}{MÉTHODE \& EXEMPLE}
\begin{frame}[t,fragile]{Lire une équation cartésienne}
Pour la droite $d:3x-2y+7=0$, une normale est $\vec n=(3,-2)$ et un directeur est $\vec u=(2,3)$.\medskip\par Le point $A(1,5)$ appartient à $d$ car $3-10+7=0$.\medskip\par Multiplier toute l’équation par un réel non nul donne la même droite : une équation n’est pas unique.
\end{frame}
\renewcommand{\slidekind}{MÉTHODE \& EXEMPLE}
\begin{frame}[t,fragile]{Construire une équation de droite}
La droite passe par $A(1,-2)$ et admet $(2,3)$ pour normale.\[2(x-1)+3(y+2)=0.\]Après développement :\[\boxed{2x+3y+4=0}.\]Contrôle : substituer les coordonnées de $A$ donne $2-6+4=0$. La normale lue dans l’équation est bien $(2,3)$.
\end{frame}
\renewcommand{\slidekind}{COMPRENDRE}
\begin{frame}[t,fragile]{Deux points donnent une direction}
Pour $A(1,2)$ et $B(4,3)$, un directeur est $\overrightarrow{AB}=(3,1)$. Une normale est donc $(1,-3)$.\medskip\par La droite $(AB)$ a pour équation\[(x-1)-3(y-2)=0,\quad\text{soit }x-3y+5=0.\]On vérifie que les deux points satisfont l’équation.
\end{frame}
\renewcommand{\slidekind}{COMPRENDRE}
\begin{frame}[t,fragile]{Parallèles et perpendiculaires}
Deux droites sont parallèles si leurs normales sont colinéaires. Elles sont perpendiculaires si leurs normales sont orthogonales.\medskip\par Les droites $2x+3y=1$ et $4x+6y=5$ sont parallèles.\medskip\par Les droites $2x+3y=1$ et $3x-2y=5$ sont perpendiculaires car $(2,3)\cdot(3,-2)=0$.
\end{frame}
\renewcommand{\slidekind}{COMPRENDRE}
\begin{frame}[t,fragile]{Projeté orthogonal : les deux conditions}
Le projeté orthogonal $H$ de $P$ sur une droite $d$ vérifie :\begin{enumerate}\item $H$ appartient à $d$ ;\item $\overrightarrow{PH}$ est colinéaire à une normale de $d$.\end{enumerate}Si $P\in d$, alors $H=P$.\medskip\par Pour calculer $H$, on intersecte $d$ avec la perpendiculaire à $d$ passant par $P$.
\end{frame}
\renewcommand{\slidekind}{MÉTHODE \& EXEMPLE}
\begin{frame}[t,fragile]{Calcul d’un projeté}
Soit $P(3,1)$ et $d:x+y-2=0$. Une normale est $(1,1)$. On pose $H=(3+t,1+t)$.\medskip\par La condition $H\in d$ donne $3+t+1+t-2=0$, donc $t=-1$.\[H=(2,0).\]Contrôle : $2+0-2=0$ et $\overrightarrow{PH}=(-1,-1)$ est bien normal à $d$.
\end{frame}
\renewcommand{\slidekind}{COMPRENDRE}
\begin{frame}[t,fragile]{Distance et point le plus proche}
Le projeté $H$ est le point de $d$ le plus proche de $P$. En effet, pour tout $M\in d$, le triangle $PHM$ est rectangle en $H$, donc\[PM^2=PH^2+HM^2\ge PH^2.\]Dans l’exemple précédent, $PH=\sqrt{(-1)^2+(-1)^2}=\sqrt2$. Cette longueur est la distance du point à la droite.
\end{frame}
\renewcommand{\slidekind}{COMPRENDRE}
\begin{frame}[t,fragile]{La médiatrice d’un segment}
Pour $A\ne B$, la médiatrice de $[AB]$ passe par le milieu $I$ et admet $\overrightarrow{AB}$ pour normale.\medskip\par Si $A(0,0)$ et $B(4,2)$, alors $I(2,1)$ et\[4(x-2)+2(y-1)=0.\]Une équation simplifiée est $2x+y-5=0$. Ses points sont équidistants de $A$ et $B$.
\end{frame}
\renewcommand{\slidekind}{MÉTHODE \& EXEMPLE}
\begin{frame}[t,fragile]{Équation d’un cercle}
Le cercle de centre $C(a,b)$ et de rayon $r>0$ est l’ensemble des $M(x,y)$ tels que\[\boxed{(x-a)^2+(y-b)^2=r^2}.\]C’est simplement l’égalité $CM^2=r^2$.\medskip\par Exemple : $(x-2)^2+(y+1)^2=9$ a pour centre $(2,-1)$ et pour rayon $3$.
\end{frame}
\renewcommand{\slidekind}{COMPRENDRE}
\begin{frame}[t,fragile]{Reconnaître un cercle par complétion}
On étudie $x^2+y^2-4x+6y+9=0$.\[x^2-4x=(x-2)^2-4,\quad y^2+6y=(y+3)^2-9.\]Ainsi\[(x-2)^2+(y+3)^2=4.\]C’est le cercle de centre $(2,-3)$ et de rayon $2$. La complétion des carrés rend la géométrie visible.
\end{frame}
\renewcommand{\slidekind}{À RETENIR}
\begin{frame}[t,fragile]{Un réel à contrôler}
Après complétion, on peut obtenir\[(x-a)^2+(y-b)^2=R.\]\begin{itemize}\item Si $R>0$ : cercle de rayon $\sqrt R$.\item Si $R=0$ : le seul point $(a,b)$.\item Si $R<0$ : aucun point réel.\end{itemize}Une équation contenant $x^2+y^2$ ne représente donc pas automatiquement un cercle.
\end{frame}
\renewcommand{\slidekind}{COMPRENDRE}
\begin{frame}[t,fragile]{Cercle de diamètre connu}
Pour $A(-1,2)$ et $B(5,2)$, le centre est le milieu $I(2,2)$ et le rayon vaut $AB/2=3$.\[\boxed{(x-2)^2+(y-2)^2=9}.\]Autre méthode : écrire $\overrightarrow{MA}\cdot\overrightarrow{MB}=0$, puis développer. Les deux méthodes produisent des équations équivalentes.
\end{frame}
\renewcommand{\slidekind}{COMPRENDRE}
\begin{frame}[t,fragile]{Intersection d’une droite et d’un cercle}
Pour le cercle $x^2+y^2=5$ et la droite $y=1$, substituer $y=1$ donne\[x^2+1=5,\qquad x^2=4.\]Les intersections sont $(-2,1)$ et $(2,1)$.\medskip\par Une droite et un cercle peuvent avoir zéro, un ou deux points communs. Un point unique correspond à une tangente.
\end{frame}
\renewcommand{\slidekind}{COMPRENDRE}
\begin{frame}[t,fragile]{Tangente à un cercle}
La tangente au cercle de centre $C$ en un point $A$ du cercle est perpendiculaire au rayon $CA$. Elle admet donc $\overrightarrow{CA}$ pour normale.\medskip\par Pour $x^2+y^2=25$ et $A(3,4)$ :\[3(x-3)+4(y-4)=0,\quad\text{soit }3x+4y=25.\]Vérifier d’abord que $A$ appartient au cercle.
\end{frame}
\renewcommand{\slidekind}{COMPRENDRE}
\begin{frame}[t,fragile]{Démontrer avec des coordonnées}
Pour étudier une configuration :\begin{enumerate}\item Choisir ou vérifier un repère orthonormé.\item Traduire les hypothèses en coordonnées ou équations.\item Effectuer les calculs exacts.\item Traduire le résultat en une conclusion géométrique.\end{enumerate}Un dessin suggère une perpendicularité, mais le produit scalaire nul la démontre.
\end{frame}
\renewcommand{\slidekind}{À RETENIR}
\begin{frame}[t,fragile]{Bilan des erreurs fréquentes}
Dans $(x-2)^2+(y+3)^2=16$, le centre est $(2,-3)$ et le rayon est $4$, pas $16$.\medskip\par Dans $ax+by+c=0$, $(a,b)$ est une \textbf{normale}, pas généralement une direction.\medskip\par Enfin, un point d’intersection doit satisfaire les \textbf{deux} équations. Toujours vérifier les coordonnées trouvées.
\end{frame}
\renewcommand{\slidekind}{COMPRENDRE}
\begin{frame}[t,fragile]{Une droite depuis un point et une normale}
La droite passe par $A(2,-1)$ et a pour normale $\vec n=(3,2)$. Un point $M(x,y)$ appartient à la droite si\[\overrightarrow{AM}\cdot\vec n=0.\] Ainsi $3(x-2)+2(y+1)=0$, soit $3x+2y-4=0$. Le point $A$ vérifie bien cette équation.
\end{frame}
\renewcommand{\slidekind}{MÉTHODE \& EXEMPLE}
\begin{frame}[t,fragile]{Un projeté résolu sans deviner}
On projette $A(3,4)$ sur $d:x+y-2=0$. Une normale est $(1,1)$. La droite perpendiculaire passant par $A$ a pour équation $y=x+1$.\medskip\par Au point d’intersection : $x+(x+1)-2=0$, donc $x=1/2$ et $y=3/2$. Le projeté est $H(1/2,3/2)$.
\end{frame}
\renewcommand{\slidekind}{COMPRENDRE}
\begin{frame}[t,fragile]{Reconnaître centre et rayon}
$x^2+y^2-4x+6y+9=0$. On complète les carrés :\[(x-2)^2-4+(y+3)^2-9+9=0.\] Donc $(x-2)^2+(y+3)^2=4$. Le centre est $(2,-3)$ et le rayon $2$. Le rayon est la racine du membre de droite, pas le membre de droite lui-même.
\end{frame}
\renewcommand{\slidekind}{À RETENIR}
\begin{frame}[t,fragile]{Contrôler une intersection}
Pour le cercle $x^2+y^2=5$ et la droite $y=x-1$, substituer donne $2x^2-2x-4=0$, soit $(x-2)(x+1)=0$.\medskip\par Les points sont $(2,1)$ et $(-1,-2)$. Vérification : chacun satisfait l’équation de la droite et celle du cercle. Une solution de l’équation en $x$ doit encore être associée à son ordonnée.
\end{frame}
\renewcommand{\slidekind}{COMPRENDRE}
\begin{frame}[t,fragile]{La normale dirige la perpendiculaire}
\begin{columns}[c,onlytextwidth]\begin{column}{.60\textwidth}\centering\begin{adjustbox}{max width=\linewidth,max height=3.7cm}\begin{tikzpicture}[>=Stealth]\draw[step=1,Line](-.3,-.3)grid(4.5,4.5);\draw[Blue,thick](-.2,2.2)--(2.2,-.2)node[right]{$d$};\draw[dashed](.5,1.5)--(3,4);\fill(3,4)circle(2pt)node[right]{$A(3,4)$};\fill(.5,1.5)circle(2pt)node[left]{$H$};\draw[->,thick,Ink](1,1)--(2,2)node[right]{$\vec n$};\end{tikzpicture}\end{adjustbox}\end{column}\begin{column}{.36\textwidth} Sur $d:x+y=2$, la normale $(1,1)$ donne la direction de $AH$. Le point $H$ appartient à $d$, tandis que $AH$ lui est perpendiculaire.\end{column}\end{columns}
\end{frame}
\renewcommand{\slidekind}{À VOUS DE JOUER}\begin{frame}[t]{Vérification des acquis 1}
Une direction de $3x+2y=7$ est…\begin{enumerate}
\item $(3,2)$
\item $(2,-3)$
\item $(2,3)$
\item $(3,-2)$
\end{enumerate}\medskip Une seule réponse est correcte. Justifier votre choix.\end{frame}
\renewcommand{\slidekind}{COMPRENDRE LA RÉPONSE}\begin{frame}[t]{Vérification 1 : résolution}
\begin{block}{Réponse 2}$(2,-3)$\end{block}\medskip $(3,2)\cdot(2,-3)=6-6=0$.\end{frame}
\renewcommand{\slidekind}{À VOUS DE JOUER}\begin{frame}[t]{Vérification des acquis 2}
La tangente à $x^2+y^2=25$ en $(5,0)$ est…\begin{enumerate}
\item $y=0$
\item $y=5$
\item $x=0$
\item $x=5$
\end{enumerate}\medskip Une seule réponse est correcte. Justifier votre choix.\end{frame}
\renewcommand{\slidekind}{COMPRENDRE LA RÉPONSE}\begin{frame}[t]{Vérification 2 : résolution}
\begin{block}{Réponse 4}$x=5$\end{block}\medskip Le rayon est horizontal, donc la tangente est verticale et passe par $(5,0)$.\end{frame}
\end{document}
